\chapter*{Appendix 12: Scenario Mu (Extraction)}

\section*{1. The Legacy Dependency (Technical \\\& Cultural)}
The physical construction and maintenance of the Sovereign Stack (hardware, nodes, solar panels) requires raw materials (silicon, copper, rare earths) that are currently locked within legacy global supply chains, fiat-denominated international trade, and ecologically destructive mining conglomerates.

\section*{2. The Parasitic Bridge (Technical \\\& Legal)}
Layer 7 utilizes automated purchasing agents (DAOs masked as legacy import/export LLCs) to source necessary raw materials from legacy markets. These entities interact with traditional supply chain APIs and use the Fiat-Crypto Bridge to settle invoices. Simultaneously, they invest surplus fiat into local circular economy initiatives (e-waste recycling, urban mining) to build sovereign material independence.

\section*{3. The Decoupling Trigger}
The bridge to global extractive supply chains is severed when the sovereign network's internal material recuperation rate (via automated recycling nodes and local closed-loop manufacturing) reaches 100\% for critical components, rendering legacy import completely obsolete.

\section*{4. Orchestration \\\& Ethical Mechanics}
This scenario exposes a critical vulnerability: the reliance on ecologically and ethically compromised legacy mining. The orchestration relies on cryptographic tracking of materials, prioritizing suppliers with the lowest collateral damage. Ethically, the parasitic bridge must be aggressively phased out; extracting materials from the legacy world is only justified if it directly accelerates the capacity for sovereign, sustainable circularity.
